\documentclass[8pt,red, compress]{beamer}
\usepackage [french]{babel}
\usepackage [utf8]{inputenc}
\usepackage{fontenc}

\title[Enfin une épi-revue en Mécanique]{{\huge Enfin une épi-revue en Mécanique !!}\\[2mm]
  Ouverte, gratuite  et éthique.}
\author{Vincent Acary, pour le comité technique}
\date{\today}

\setbeamertemplate 
{footline} 
{\quad\hfill\strut\insertsection\quad--\quad\insertframenumber/\inserttotalframenumber\strut\quad\quad} 

\newcommand{\tr}[1]{\textcolor{red}{#1}}
\usepackage{smartdiagram}


\usepackage{hyperref}
\begin{document}
{
  \begin{frame}[plain,noframenumbering]
    \begin{center}
     \includegraphics[width=1.0\textwidth]{BandeauWeb.pdf}
   \end{center}
    \titlepage
  \end{frame}
}
\begin{frame}
  \begin{block}{Une longue démarche}
  \begin{itemize}
  \item Sep. 2015. Première discussion
  \item Jun. 2017. Prise de contact au CSMA
  \item Jan. 2018. AG Mécamat Aussois "Pour une épi-revue en Mécanique
  \item Aujourd'hui, un premier papier accepté et 15 soumissions en relecture.
    \begin{center}
    \includegraphics[width=0.6\textwidth]{first_paper.png}
  \end{center}
  \end{itemize}
\end{block}
\begin{center}
\Huge{\tr{Un grand merci à Mécamat \\
    pour son soutien de la première heure}}
\end{center}
\end{frame}


\end{document}

%%% Local Variables: 
%%% mode: latex
%%% TeX-master: t
%%% End: 